\documentclass[11pt]{article}
\usepackage{../common/polystyle}
\usepackage{xr}
\externaldocument[P-]{poly_pacbayes}
\graphicspath{{figures/}}
\renewcommand{\polytitle}{PAC-Bayesian Learning -- Solutions to the exercises}
\newcommand{\statement}[1]{\begin{tcolorbox}[enhanced,breakable,colback=black!3,colframe=black!25,boxrule=0.4pt,sharp corners,fontupper=\itshape\small]#1\end{tcolorbox}}
\newcommand{\exo}[1]{\Needspace{6\baselineskip}\subsection*{Exercise #1}\addcontentsline{toc}{subsection}{Exercise #1}}
\newcommand{\sol}{\noindent\textbf{Solution.}\ }

\begin{document}
\sloppy
\begin{center}
{\LARGE\bfseries PAC-Bayesian Learning}\\[0.2cm]
{\Large Detailed solutions to the exercises of the lecture notes}\\[0.3cm]
{\large Master MALIA -- Ensemble Methods}\\[0.4cm]
Guillaume Metzler\\
Institut de Communication (ICOM), Université Lumière Lyon 2\\
Laboratoire ERIC UR 3083, Lyon, France
\end{center}

\bigskip
\noindent This document gathers detailed solutions to the exercises that close each section of the lecture notes
\emph{PAC-Bayesian Learning: A Theoretical Framework for Ensemble Methods}. Exercises, figures, equations and results
are numbered as in the notes. The numerical values were computed with \texttt{pbtools.py}; the short scripts used
for the simulation exercises are those of the notebooks. We advise to try each exercise before reading its
solution: most of them take a few lines once the right tool of the notes has been identified, and the remarks
that follow the solutions point to what the exercise is meant to teach.

\tableofcontents
\bigskip

%% ======================================================================
\section*{Section 1 -- Why PAC-Bayes for ensemble methods?}
\addcontentsline{toc}{section}{Section 1 -- Why PAC-Bayes for ensemble methods?}

\exo{1.1}
\statement{Write the predictors of bagging, of AdaBoost and of a kernel machine with random Fourier features in the
form \eqref{P-eq:mv}. In each case, what is the set of voters $\Hcal$, and what is the distribution $Q$?}

\sol In all three cases we look for a set of voters $\Hcal$ and a probability distribution $Q$ on $\Hcal$ such that the
predictor equals $\sign(\E_{h\sim Q}h(x))$.

\emph{Bagging and random forests.} The voters are the $n$ trees $h_1,\dots,h_n$ learned on the bootstrap samples, so
$\Hcal=\{h_1,\dots,h_n\}$, and $Q$ is the uniform distribution $Q(h_t)=\frac1n$. Indeed
$\sign\big(\frac1n\sum_th_t(x)\big)=\sign(\E_{h\sim Q}h(x))$.

\emph{AdaBoost.} The voters are the weak learners $h_1,\dots,h_T$, and $Q(h_t)=\alpha_t/\sum_s\alpha_s$. The weights are
non-negative because each weak learner has a weighted error $\varepsilon_t<\frac12$, so that
$\alpha_t=\frac12\ln\frac{1-\varepsilon_t}{\varepsilon_t}>0$, and dividing the score $\sum_t\alpha_th_t(x)$ by the positive
constant $\sum_s\alpha_s$ does not change its sign. If some $\alpha_t$ were negative, we would replace $h_t$ by $-h_t$ and
$\alpha_t$ by $|\alpha_t|$: a vote with signed weights is always a vote with non-negative weights over the set
$\Hcal\cup(-\Hcal)$.

\emph{Random Fourier features.} The kernel machine is $\sign\big(\sum_{k=1}^K\alpha_k\cos(\omega_k^\top x+b_k)\big)$. The voters are
the functions $h_{\omega,b,s}(x)=s\cos(\omega^\top x+b)$ with $s\in\{-1,+1\}$ (which take values in $[-1,1]$ rather than
$\{-1,+1\}$; for binary voters take $s\,\sign\cos(\omega^\top x+b)$ as in \cref{P-fig:rff}), and $Q$ puts the weight
$|\alpha_k|/\sum_j|\alpha_j|$ on the voter $(\omega_k,b_k,\sign\alpha_k)$. The set of voters is infinite, indexed by
$(\omega,b,s)$, and the distribution from which the frequencies are drawn (the Fourier transform of the kernel,
times the uniform distribution of $b$ on $[0,2\pi]$) plays the role of a \emph{prior}: this is the viewpoint of
\citet{letarte2019pseudo} discussed in \cref{P-sec:linear}.

\emph{Remark.} The exercise shows that the question ``which predictor?'' reduces, for all these methods, to the
question ``which distribution $Q$ over voters?'', which is precisely the object of PAC-Bayes.

\exo{1.2}
\statement{Show that the empirical Rademacher complexity of $\mathrm{conv}(\Hcal)$ equals the one of $\Hcal$.
\emph{Hint:} for a fixed vector of signs $\sigma$, the function $f\mapsto\frac1m\sum_i\sigma_if(x_i)$ is linear in $f$.}

\sol Recall that $\widehat{\mathfrak R}_S(\mathcal F)=\E_\sigma\big[\sup_{f\in\mathcal F}\frac1m\sum_i\sigma_if(x_i)\big]$, where the
$\sigma_i$ are independent uniform signs. Since $\Hcal\subseteq\mathrm{conv}(\Hcal)$, the supremum over the convex hull is at
least the supremum over $\Hcal$, so $\widehat{\mathfrak R}_S(\Hcal)\leq\widehat{\mathfrak R}_S(\mathrm{conv}(\Hcal))$.

For the converse, fix $\sigma$ and let $L_\sigma(f)=\frac1m\sum_i\sigma_if(x_i)$, which is linear in $f$. Any
$f\in\mathrm{conv}(\Hcal)$ can be written $f=\sum_{j=1}^kQ_jh_j$ with $h_j\in\Hcal$, $Q_j\geq0$ and $\sum_jQ_j=1$. By linearity,
\[
L_\sigma(f)=\sum_jQ_jL_\sigma(h_j)\leq\max_jL_\sigma(h_j)\leq\sup_{h\in\Hcal}L_\sigma(h).
\]
Taking the supremum over $f$ and then the expectation over $\sigma$ gives
$\widehat{\mathfrak R}_S(\mathrm{conv}(\Hcal))\leq\widehat{\mathfrak R}_S(\Hcal)$, hence equality. (For infinite convex combinations
the same argument applies with integrals.)

\emph{Remark.} Averaging voters does not increase the capacity measured by the Rademacher complexity. This is the
key to the margin bounds of \citet{schapire1998boosting} and \citet{koltchinskii2002empirical}, which do not depend on
the number of voters in the vote.

\exo{1.3}
\statement{In \cref{P-fig:adaboost}, the training error reaches zero after a few dozen rounds. Using the right panel,
explain why a bound based on the empirical margin distribution can keep decreasing afterwards.}

\sol A margin bound (\cref{P-sec:margins}) has the form
$\RD(\BQ)\leq\P_S(M_Q\leq\theta)+O\big(\sqrt{\ln|\Hcal|\ln m/(m\theta^2)}\big)$ for every $\theta>0$. Its empirical term is not the
training error, which is $\P_S(M_Q\leq0)$, but the proportion of training examples whose normalized margin is below
$\theta$. Once the training error is zero, AdaBoost keeps increasing the margins of the examples: the right panel of
\cref{P-fig:adaboost} shows that the empirical distribution function of the margins moves to the right between $T=50$ and
$T=400$ rounds. For a fixed $\theta$, the quantity $\P_S(M_Q\leq\theta)$ therefore keeps decreasing; equivalently, we can
use a larger $\theta$ for the same empirical term, which decreases the complexity term in $1/\theta^2$. The bound thus
keeps improving although the training error no longer moves, which is consistent with the stable (and even slightly
decreasing) test error of the middle panel.

%% ======================================================================
\section*{Section 2 -- Mathematical toolbox}
\addcontentsline{toc}{section}{Section 2 -- Mathematical toolbox}

\exo{2.1}
\statement{Compute $\kl(0.1\|0.2)$ and $\kl(0.2\|0.1)$. Is the binary kl symmetric?}

\sol By definition $\kl(q\|p)=q\ln\frac qp+(1-q)\ln\frac{1-q}{1-p}$. Hence
\begin{align*}
\kl(0.1\|0.2)&=0.1\ln\tfrac12+0.9\ln\tfrac98=-0.0693+0.1060=0.0367,\\
\kl(0.2\|0.1)&=0.2\ln2+0.8\ln\tfrac89=0.1386-0.0942=0.0444.
\end{align*}
The two values differ: the binary kl is not symmetric. The interpretation of \cref{P-chap:tools} helps: $\kl(q\|p)$ measures
the surprise of observing a frequency $q$ when the true rate is $p$. Observing $0.1$ errors when the true rate is
$0.2$ is less surprising than observing $0.2$ when the true rate is $0.1$, because a Bernoulli variable of mean $0.2$
fluctuates more than one of mean $0.1$.

\exo{2.2}
\statement{Show that the Gibbs distribution $Q_\phi$ of Lemma~\ref{P-lem:dv} is the only posterior for which the change of
measure is an equality.}

\sol The proof of Lemma~\ref{P-lem:dv} shows that, for every $Q$ with $\KL(Q\|P)<\infty$,
\[
\KL(Q\|P)+\ln\E_Pe^\phi-\E_Q\phi=\KL(Q\|Q_\phi).
\]
The change of measure is therefore an equality iff $\KL(Q\|Q_\phi)=0$. By Gibbs' inequality (the KL is non-negative and
vanishes only for equal distributions, see TD Exercise~1), this happens iff $Q=Q_\phi$. If $\KL(Q\|P)=+\infty$, the
right-hand side of \eqref{P-eq:dv} is infinite while the left-hand side is finite (when $\E_Q\phi$ is), so there is no
equality either. Hence $Q_\phi$ is the unique posterior achieving equality, and the supremum in
\eqref{P-eq:gibbs_variational} is attained only at $Q_\phi$.

\exo{2.3}
\statement{Using the Occam bound in Pinsker form, how many examples are needed to certify a risk below $0.15$ for a
classifier with empirical risk $0.1$ chosen among $10^6$ candidates, with $\delta=0.05$?}

\sol With the uniform prior on $|\Hcal|=10^6$ classifiers, $\ln\frac1{P(h)}=\ln10^6=13.82$, and $\ln\frac1\delta=\ln20=3.00$. The
Pinsker form of Theorem~\ref{P-thm:occam} gives
$\RD(h)\leq0.1+\sqrt{\frac{13.82+3.00}{2m}}$. The certificate is below $0.15$ iff
\[
\sqrt{\frac{16.81}{2m}}\leq0.05\iff m\geq\frac{16.81}{2\times0.0025}=3362.2,
\]
so $m=3363$ examples suffice. \emph{Remarks.} With the kl form, far fewer examples are needed: about $1550$, since the kl
is much tighter for small risks. And the dependence on the number of candidates is only logarithmic: with $10^{12}$
candidates, the numerator becomes $27.63+3.00$ and we need $6127$ examples, less than twice as many.

\exo{2.4}
\statement{Using Pinsker's inequality, show that the kl interval of \cref{P-fig:concentration} is never wider than the
Hoeffding interval. Using notebook~1, determine for which empirical risks it is at least twice narrower when $m=1000$.}

\sol Let $\varepsilon=\ln(1/\delta)/m$. The Hoeffding interval is $[\hat q,\hat q+\sqrt{\varepsilon/2}]$ and the kl interval is
$[\hat q,\klup(\hat q,\varepsilon)]$. If $p$ belongs to the kl interval, then $\kl(\hat q\|p)\leq\varepsilon$, and Pinsker's inequality
gives $2(p-\hat q)^2\leq\kl(\hat q\|p)\leq\varepsilon$, i.e.\ $p\leq\hat q+\sqrt{\varepsilon/2}$. Hence
$\klup(\hat q,\varepsilon)\leq\hat q+\sqrt{\varepsilon/2}$: the kl interval is always contained in the Hoeffding interval.

Numerically, for $m=1000$ and $\delta=0.05$, the Hoeffding width is $\sqrt{\ln20/2000}=0.0387$. Computing
$\klup(\hat q,\varepsilon)-\hat q$ on a grid of $\hat q$, the ratio of the two widths is $3.97$ at $\hat q=0.01$, still $2$ at
$\hat q\approx0.054$, and decreases to $1.00$ at $\hat q=0.5$. The kl interval is thus at least twice narrower for all
empirical risks below about $5.4\%$. Pinsker's inequality is tight around $\frac12$, which is why the two intervals
coincide there.

%% ======================================================================
\section*{Section 3 -- PAC-Bayesian generalization bounds}
\addcontentsline{toc}{section}{Section 3 -- PAC-Bayesian generalization bounds}

\exo{3.1}
\statement{Starting from Theorem~\ref{P-thm:general} with $\Delta(q,p)=2(q-p)^2$, prove McAllester's bound directly,
without going through the kl.}

\sol The function $\Delta(q,p)=2(q-p)^2$ is jointly convex (it is a convex function of the linear form $q-p$), so
Theorem~\ref{P-thm:general} applies. It remains to bound $\mathcal I_\Delta(m)=\E_{S'}\E_{h\sim P}e^{2m(\widehat R_{S'}(h)-\RD(h))^2}$.
For a fixed $h$, by Pinsker's inequality, $2(\hat p-p)^2\leq\kl(\hat p\|p)$, so
$\E_{S'}e^{2m(\hat p-p)^2}\leq\E_{S'}e^{m\kl(\hat p\|p)}\leq2\sqrt m$ by Maurer's lemma (Lemma~\ref{P-lem:chernoff}, for $m\geq8$).
Hence $\mathcal I_\Delta(m)\leq2\sqrt m$ and, with probability at least $1-\delta$, for all $Q$,
\[
2\big(\RS(\GQ)-\RD(\GQ)\big)^2\leq\frac{\KL(Q\|P)+\ln\frac{2\sqrt m}{\delta}}{m},
\]
hence
\[
\RD(\GQ)\leq\RS(\GQ)+\sqrt{\frac{\KL(Q\|P)+\ln\frac{2\sqrt m}{\delta}}{2m}}.
\]
\emph{Remark.} One can also bound $\E e^{2m(\hat p-p)^2}$ without the kl, using the sub-Gaussian tail
$\P(|\hat p-p|\geq t)\leq2e^{-2mt^2}$ and the formula $\E e^{Z}=1+\int_0^\infty e^z\P(Z>z)\,dz$; this gives a bound of order
$m$ instead of $2\sqrt m$, i.e.\ a slightly worse constant ($\ln(2m)$ instead of $\ln(2\sqrt m)$). This is the form of the
original bounds of McAllester.

\exo{3.2}
\statement{For $m=500$, $\delta=0.05$, $\RS(\GQ)=0.05$ and $\KL(Q\|P)\in\{0,5,20\}$, compute McAllester's and Seeger's bounds,
and compare with \cref{P-tab:bound_anatomy}.}

\sol The confidence term is $\ln\frac{2\sqrt{500}}{0.05}=\ln894.4=6.80$. The complexity per example
$\varepsilon=\frac{\KL+6.80}{500}$ and the two bounds are:
\begin{center}\small
\begin{tabular}{cccc}
\toprule
$\KL(Q\|P)$ & $\varepsilon$ & McAllester $0.05+\sqrt{\varepsilon/2}$ & Seeger $\klup(0.05,\varepsilon)$ \\
\midrule
$0$ & $0.0136$ & $0.132$ & $0.094$ \\
$5$ & $0.0236$ & $0.159$ & $0.112$ \\
$20$ & $0.0536$ & $0.214$ & $0.154$ \\
\bottomrule
\end{tabular}
\end{center}
Compared with \cref{P-tab:bound_anatomy} ($m=1000$, $\RS=0.12$), the gap between the two bounds is larger here: the
complexity term of Seeger's bound is about $30\%$ smaller than McAllester's when $\RS=0.12$, but about $40\%$ smaller when
$\RS=0.05$. This is the fast rate of the kl at work: the smaller the empirical risk, the more the square root of
Pinsker's inequality is wasteful. Halving $m$ roughly doubles $\varepsilon$, which costs Seeger's bound less than
McAllester's for the same reason.

\exo{3.3}
\statement{Using the linear form of Catoni's bound given after Corollary~\ref{P-cor:catoni}, find the value of $C$ that
minimizes it when $\RS(\GQ)=0.12$, $\KL(Q\|P)=5$, $m=1000$ and $\delta=0.05$, and compare the result with
\cref{P-tab:bound_anatomy}.}

\sol Let $A=\frac{\KL+\ln(1/\delta)}m=\frac{5+3.00}{1000}=0.0080$ and $f(C)=\frac{C\RS+A}{1-e^{-C}}$. Setting
$f'(C)=0$ gives $\RS(1-e^{-C})=(C\RS+A)e^{-C}$, i.e.\ $\RS(e^C-1-C)=A$. With $\RS=0.12$, $e^C-1-C=0.0667$; since
$e^C-1-C\approx\frac{C^2}2+\frac{C^3}6$, a first approximation is $C\approx\sqrt{2A/\RS}=0.365$, and solving numerically gives
$C^\star=0.344$ and $f(C^\star)=0.169$. The exact (non-linear) Catoni bound is minimized at $C=0.371$, where it equals
$0.165$. In \cref{P-tab:bound_anatomy}, Catoni's bound optimized on a grid of $30$ values of $C$ (with a union bound, hence
$\ln30$ added to the numerator) gives $0.166$ for $\KL=5$ and Seeger's bound $0.168$: the linear relaxation loses
only $0.004$ with respect to the exact bound, and the grid costs very little. \emph{Warning:} as for $\lambda$ in the
Hoeffding bound, $C^\star$ depends on $\RS$ and $\KL$, hence on the data; using it directly requires a grid and a union bound.

\exo{3.4}
\statement{Reproduce the experiment of \cref{P-fig:uniformity} for $m\in\{100,1000,10000\}$. Show that the failure frequency
of the naive bound does not go to zero when $m$ grows, and explain why by comparing $\sqrt{\ln(1/\delta)/(2m)}$ with
the typical deviation of the minimum of $1000$ empirical risks.}

\sol We draw $1000$ classifiers of true risk $0.5$ (their empirical risks are independent $\mathrm{Bin}(m,\frac12)/m$),
keep the smallest empirical risk $\hat r_{\min}$, and check whether the naive bound $\hat r_{\min}+\sqrt{\ln(1/\delta)/(2m)}$ is
below $0.5$. Over $2000$ repetitions:
\begin{center}\small
\begin{tabular}{cccccc}
\toprule
$m$ & mean of $0.5-\hat r_{\min}$ & naive width & ratio & naive bound wrong & Occam wrong \\
\midrule
$100$ & $0.160$ & $0.122$ & $1.31$ & $99.8\%$ & $0.3\%$ \\
$1000$ & $0.051$ & $0.039$ & $1.32$ & $99.9\%$ & $0.5\%$ \\
$10000$ & $0.016$ & $0.012$ & $1.32$ & $100\%$ & $0.6\%$ \\
\bottomrule
\end{tabular}
\end{center}
The failure frequency does not decrease. The reason is that both quantities scale as $1/\sqrt m$. Each empirical risk
is approximately Gaussian with standard deviation $\frac{1}{2\sqrt m}$, and the minimum of $n$ independent Gaussian
variables lies about $\sqrt{2\ln n}$ standard deviations below their mean. The selection bias is therefore of order
$\frac{\sqrt{2\ln n}}{2\sqrt m}=\frac{1.86}{\sqrt m}$ for $n=1000$ (in practice $\approx1.6/\sqrt m$ at these sample sizes), while
the naive width is $\frac{\sqrt{\ln(1/\delta)/2}}{\sqrt m}=\frac{1.22}{\sqrt m}$. Their ratio does not depend on $m$: the naive bound
is always too narrow by the same factor. The Occam bound has width $\frac{\sqrt{(\ln n+\ln(1/\delta))/2}}{\sqrt m}=\frac{2.22}{\sqrt m}$,
which exceeds the selection bias: the price $\ln n$ is exactly what is needed to absorb the selection. More data do not
cure selection bias; only accounting for the selection does.

%% ======================================================================
\section*{Section 4 -- Posteriors, priors and the computation of the bounds}
\addcontentsline{toc}{section}{Section 4 -- Posteriors, priors and the computation of the bounds}

\exo{4.1}
\statement{For a finite class and a uniform prior, show that $\KL(Q_\lambda\|P)\to\ln\frac{|\Hcal|}{k}$ when $\lambda\to\infty$, where $k$
is the number of empirical risk minimizers.}

\sol Let $n=|\Hcal|$, $r^\star=\min_h\RS(h)$ and $\mathcal M=\{h:\RS(h)=r^\star\}$, with $|\mathcal M|=k$. With the uniform prior,
\[
Q_\lambda(h)=\frac{e^{-\lambda m\RS(h)}}{\sum_{h'}e^{-\lambda m\RS(h')}}=\frac{e^{-\lambda m(\RS(h)-r^\star)}}{k+\sum_{h'\notin\mathcal M}e^{-\lambda m(\RS(h')-r^\star)}}.
\]
For $h'\notin\mathcal M$, $\RS(h')-r^\star>0$, so the exponentials tend to $0$ when $\lambda\to\infty$. Hence $Q_\lambda(h)\to\frac1k$ for
$h\in\mathcal M$ and $Q_\lambda(h)\to0$ otherwise: $Q_\lambda$ converges to the uniform distribution $U_{\mathcal M}$ on the
minimizers. The KL is a continuous function of $Q$ on the simplex (with $0\ln0=0$) when $P$ has full support, so
$\KL(Q_\lambda\|P)\to\KL(U_{\mathcal M}\|P)=\sum_{h\in\mathcal M}\frac1k\ln\frac{1/k}{1/n}=\ln\frac nk$ (TD Exercise~1). With a unique minimizer we
recover the Occam price $\ln n$; when several classifiers tie, the Gibbs posterior spreads over them and pays less.

\exo{4.2}
\statement{Explain why using the same sample to learn the centre of a Gaussian prior and to compute the empirical term
of the bound is not allowed. What happens if we do it anyway?}

\sol The only probabilistic step of the proof of Theorem~\ref{P-thm:general} is the third one, where we write
$\E_{S'}\E_{h\sim P}e^{m\Delta(\widehat R_{S'}(h),\RD(h))}=\E_{h\sim P}\E_{S'}e^{m\Delta(\cdots)}$ and bound the inner expectation for each
\emph{fixed} $h$. This requires $P$ not to depend on $S'$: if $P=P_S$ is centred on a classifier learned on $S$, the
voters drawn from $P_S$ are correlated with the sample, their empirical risks are biased downwards, and the inner
expectation can no longer be controlled by Maurer's lemma. The theorem simply does not apply.

If we do it anyway, the ``certificate'' can be badly wrong. Take the extreme case where the posterior equals the prior,
both centred on the learned classifier: the KL is zero, and the bound becomes a concentration inequality for a
classifier that was selected with the data---exactly the naive bound of \cref{P-fig:uniformity}, which fails almost surely
when the selection is strong. With a flexible model (for instance a neural network that interpolates the data),
the empirical risk is close to zero, the KL is zero, and the certificate would claim a near-zero risk regardless of
the true one. The correct procedure is to split $S$: learn the prior on $S_1$ and compute the bound on $S_2$
(Proposition~\ref{P-prop:ddprior}, \cref{P-fig:ddprior}).

\exo{4.3}
\statement{In the Monte Carlo certificate, how should $\delta$ and $\delta'$ be split to minimize the final bound when $N=100$?
when $N=10^4$?}

\sol The final certificate holds with probability $1-\delta_{\text{tot}}$ with $\delta_{\text{tot}}=\delta+\delta'$; we fix
$\delta_{\text{tot}}=0.05$ and choose the fraction $\alpha=\delta'/\delta_{\text{tot}}$ given to the Monte Carlo step. Increasing $\alpha$
makes the Monte Carlo inversion $\klup(\hat r_N,\ln(2/\delta')/N)$ tighter but the PAC-Bayesian bound looser (it uses
$\delta=(1-\alpha)\delta_{\text{tot}}$). The two terms enter only through logarithms, so the optimum is flat. On the example of
\cref{P-fig:montecarlo} (\texttt{breast cancer}, $\KL=32$, $m=398$), a grid search over $\alpha$ gives the following. With
$N=100$ samples, the Monte Carlo term is the main source of looseness, since $\ln(2/\delta')/N$ is large; it deserves most
of the budget, and the optimum is $\alpha\approx0.86$ ($\delta'=0.043$, $\delta=0.007$), with a certificate of $0.380$ against $0.386$
for the even split. With $N=10^4$ samples, the Monte Carlo term is small and the PAC-Bayesian term dominates; the
optimum moves to $\alpha\approx0.3$, and the certificate ($0.2591$) is practically the same as with the even split ($0.2594$).
In practice the even split is a sensible default; the gain from optimizing is small, and the split must anyway be
chosen before computing the bound (or on a grid with a union bound). It is more effective to increase $N$.

%% ======================================================================
\section*{Section 5 -- From the Gibbs classifier to the majority vote}
\addcontentsline{toc}{section}{Section 5 -- From the Gibbs classifier to the majority vote}

\exo{5.1}
\statement{Give an example of three voters and a distribution on three points for which $\RD(\BQ)=0$ while $\RD(\GQ)=\frac13$.
Compute the three oracle bounds.}

\sol Let $\D$ be uniform on three points $z_1,z_2,z_3$, and let voter $h_i$ err on $z_i$ only; $Q$ is uniform on
$\{h_1,h_2,h_3\}$. On each point exactly one voter out of three errs, so $W_Q(z_j)=\frac13$ for all $j$. The vote is correct
everywhere ($W_Q<\frac12$): $\RD(\BQ)=0$, while $\RD(\GQ)=\E W_Q=\frac13$. The second-order quantities are
$e_Q=\E W_Q^2=\frac19$ and $d_Q=\E[2W_Q(1-W_Q)]=2\cdot\frac13\cdot\frac23=\frac49$. Hence:
\[
2\RD(\GQ)=\tfrac23,\qquad 4e_Q=\tfrac49,\qquad
\mathcal C_Q=1-\frac{(1-2/3)^2}{1-8/9}=1-\frac{1/9}{1/9}=0
\]
for the first-order bound, the tandem bound and the C-bound respectively.
The C-bound is exact here: $W_Q$ is constant, its variance is zero, and Cantelli's inequality says that a variable with
zero variance equal to $\frac13$ is never above $\frac12$. The example is a small Condorcet jury with perfectly
``anti-correlated'' errors: the disagreement ($\frac49$) exceeds the Gibbs risk ($\frac13$), so Proposition~\ref{P-prop:compare}
predicts the order C-bound $\leq$ tandem $\leq$ first order, as observed.

\exo{5.2}
\statement{Show that the C-bound is always smaller than $1$ when $\RD(\GQ)<\frac12$, and that it equals $4\RD(\GQ)(1-\RD(\GQ))$ when
$\dD(Q)=0$. Compare with the first-order bound.}

\sol Write $R=\RD(\GQ)$ and $d=\dD(Q)$. By Proposition~\ref{P-prop:decomp}, $0\leq d\leq2R(1-R)<\frac12$, so $1-2d>0$. When $R<\frac12$,
$(1-2R)^2>0$, hence $\frac{(1-2R)^2}{1-2d}>0$ and $\mathcal C_Q<1$. (In fact $1-2d\geq1-4R+4R^2=(1-2R)^2$, so the ratio is at most $1$
and $\mathcal C_Q\geq0$.) When $d=0$, $\mathcal C_Q=1-(1-2R)^2=4R-4R^2=4R(1-R)$. Compared with the first-order bound $2R$:
$4R(1-R)\geq2R\iff2(1-R)\geq1\iff R\leq\frac12$. So when the voters never disagree, the C-bound is \emph{worse} than the
first-order bound, by a factor $2(1-R)$ which is close to $2$ for small $R$. This is consistent with
Proposition~\ref{P-prop:compare} ($d=0<R$): without diversity, second-order information can only hurt, since all voters
predict the same thing and the vote is just one of them.

\exo{5.3}
\statement{$(\star)$ Prove that the minimum over $\mu$ of the Chebyshev--Cantelli family equals the C-bound (TD, Exercise~5).}

\sol Let $R=\E W_Q<\frac12$, $e=\E W_Q^2$, $v=e-R^2=\Var W_Q$, and $f(\mu)=\frac{e-2\mu R+\mu^2}{(1/2-\mu)^2}$ for $\mu<\frac12$. The
numerator equals $\E(W_Q-\mu)^2=v+(R-\mu)^2$. Differentiating,
\[
f'(\mu)=\frac{2(\mu-R)(\tfrac12-\mu)+2\big(v+(R-\mu)^2\big)}{(1/2-\mu)^3}=\frac{2\big[v-(R-\mu)(\tfrac12-R)\big]}{(1/2-\mu)^3},
\]
using $(\mu-R)(\frac12-\mu)+(R-\mu)^2=(R-\mu)(R-\frac12)$. The bracket is increasing in $\mu$ and vanishes at
$\mu^\star=R-\frac{v}{1/2-R}=\frac{R/2-e}{1/2-R}$; $f$ decreases before $\mu^\star$ and increases after, so $\mu^\star$ is the minimum
(and $\mu^\star<R<\frac12$ is admissible). At $\mu^\star$, $R-\mu^\star=\frac{v}{a}$ with $a=\frac12-R$, and
$\frac12-\mu^\star=a+\frac va=\frac{a^2+v}{a}$. Hence
\[
f(\mu^\star)=\frac{v+v^2/a^2}{(a^2+v)^2/a^2}=\frac{v(a^2+v)}{(a^2+v)^2}=\frac{v}{v+a^2}.
\]
Finally, with the margin $M_Q=1-2W_Q$: $v=\frac14\Var M_Q$ and $a^2=\frac14(\E M_Q)^2$, so
$f(\mu^\star)=\frac{\Var M_Q}{\Var M_Q+(\E M_Q)^2}=1-\frac{(\E M_Q)^2}{\E M_Q^2}=\mathcal C_Q$, the C-bound \eqref{P-eq:cbound}. The tandem
bound is $f(0)$, hence always at least the C-bound, which proves the first claim of Proposition~\ref{P-prop:compare}.

\exo{5.4}
\statement{In Condorcet's model with $n$ independent voters of error $p$, compute $\eD(Q)$ and $\dD(Q)$, and show that the
C-bound decreases like $1/n$ while the tandem bound tends to $4p^2$.}

\sol Since $nW_Q\sim\mathrm{Bin}(n,p)$, $\E W_Q=p$ and $\Var W_Q=\frac{p(1-p)}n$. Hence
$\eD(Q)=\E W_Q^2=p^2+\frac{p(1-p)}n$ and, by Proposition~\ref{P-prop:decomp}, $\dD(Q)=2(p-\eD(Q))=2p(1-p)\big(1-\frac1n\big)$. The tandem
bound is $4\eD(Q)=4p^2+\frac{4p(1-p)}n\to4p^2$: it never goes to zero, even though the risk of the vote does (for $p=0.3$
it tends to $0.36$). For the C-bound,
\begin{align*}
1-2\dD(Q)&=1-4p(1-p)+\frac{4p(1-p)}n=(1-2p)^2+\frac{4p(1-p)}n,\\
\mathcal C_Q&=1-\frac{(1-2p)^2}{1-2\dD(Q)}=\frac{4p(1-p)/n}{(1-2p)^2+4p(1-p)/n}.
\end{align*}
When $n\to\infty$, $\mathcal C_Q\sim\frac{4p(1-p)}{(1-2p)^2}\cdot\frac1n$. For $p=0.3$ the equivalent is $\frac{0.84}{0.16\,n}=\frac{5.25}n$; the exact values are $0.174$ for $n=25$ and $0.049$ for
$n=101$ (\cref{P-fig:condorcet}). The true risk decreases exponentially fast (by Hoeffding's inequality,
$\RD(\BQ)\leq e^{-2n(1/2-p)^2}$), so even the C-bound is pessimistic, but it captures the qualitative effect of independence.
The tandem bound misses it because it applies Markov's inequality to $W_Q^2$, i.e.\ it measures the deviations of
$W_Q$ from $0$: since $W_Q$ concentrates around $p>0$, $\E W_Q^2\to p^2$ does not vanish. The C-bound measures the deviations
from the well-chosen centre $\mu^\star$ (Exercise~5.3), and only pays for the variance $\frac{p(1-p)}n$, which vanishes.

%% ======================================================================
\section*{Section 6 -- Self-bounding learning}
\addcontentsline{toc}{section}{Section 6 -- Self-bounding learning}

\exo{6.1}
\statement{Reproduce the experiment of \cref{P-fig:uniformity} (draw $1000$ random classifiers and $m=100$ labels). How often
is the ``fixed-$h$'' Hoeffding bound violated? And the Occam bound? What happens to the Occam bound with $10^6$
classifiers?}

\sol See Exercise~3.4 for the code and the explanation. With $1000$ classifiers and $m=100$, the best empirical risk is
$0.34$ on average; the fixed-$h$ Hoeffding bound ($+0.122$) is violated in about $99.8\%$ of the runs, while the Occam
bound ($+\sqrt{(\ln1000+\ln20)/200}=+0.222$) is violated in about $0.3\%$ of them, below the nominal $5\%$. With $10^6$
classifiers, the best empirical risk drops to $0.26$ on average (a stronger selection), but the Occam price grows only
to $\sqrt{(\ln10^6+\ln20)/200}=0.290$: the certificate, about $0.55$ on average, is never violated in $200$ runs. It is
vacuous, which is the honest answer: with $100$ examples one cannot distinguish a million random classifiers from each
other. The logarithmic price keeps the bound valid for any number of candidates, at the cost of informativeness.

\exo{6.2}
\statement{Show that the value of $\lambda$ minimizing the Hoeffding-type bound \eqref{P-eq:catoni_hoeffding} depends on
$\KL(Q\|P)$, and explain why using it directly is forbidden, whereas using $\lambda^\star$ in the PAC-Bayes-$\lambda$ bound is allowed.}

\sol The bound is $\RS(\GQ)+\frac{A}{\lambda}+\frac\lambda8$ with $A=\frac{\KL(Q\|P)+\ln(1/\delta)}m$; its derivative in $\lambda$ vanishes
at $\lambda^\star=\sqrt{8A}=\sqrt{8(\KL(Q\|P)+\ln\frac1\delta)/m}$, which depends on $\KL(Q\|P)$, hence on the posterior and on the data.
In Corollary~\ref{P-cor:hoeffding_pb}, $\lambda$ is a parameter of the comparison function $\Delta(q,p)=\lambda(p-q)$: each value of
$\lambda$ defines a different event of probability at least $1-\delta$, and the corollary guarantees only that the event
associated with a $\lambda$ \emph{fixed in advance} holds. Choosing $\lambda$ after seeing the data means picking one event
among infinitely many, and the probability that the picked one fails is not controlled (the situation of
\cref{P-fig:uniformity}). It must be done on a finite grid with a union bound (Proposition~\ref{P-prop:union}).

In Corollary~\ref{P-cor:lambda}, on the contrary, the probabilistic event is the single event of Seeger's bound, which does
not involve $\lambda$; the PAC-Bayes-$\lambda$ inequality is derived from it by deterministic algebra valid for every
$\lambda\in(0,2)$ simultaneously. Plugging the data-dependent $\lambda^\star$ is therefore legitimate: the rule is ``free if it
indexes a family of inequalities all implied by the same event, paid if it changes the event''
(\cref{P-sec:sb_rules}).

\exo{6.3}
\statement{Using Proposition~\ref{P-prop:klgrad}, compute $\partial p/\partial q$ at $\hat q=0.1$, $\varepsilon=0.05$ (with $p\approx0.220$). Which
ingredient of the certificate is it more profitable to decrease: the empirical risk by $0.01$, or the KL by $5$ nats with
$m=1000$?}

\sol With $p=0.2201$, $q=0.1$: $\frac{p(1-p)}{p-q}=\frac{0.2201\times0.7799}{0.1201}=1.429$ and
$\ln\frac{p(1-q)}{q(1-p)}=\ln\frac{0.1981}{0.0780}=\ln2.540=0.932$. Hence
\[
\frac{\partial p}{\partial q}=1.429\times0.932=1.33,\qquad\frac{\partial p}{\partial\varepsilon}=1.43.
\]
Decreasing the empirical risk by $0.01$ decreases the certificate by about $1.33\times0.01=0.0133$ (the exact value is
$0.0135$). Decreasing the KL by $5$ nats decreases $\varepsilon$ by $5/1000=0.005$ and the certificate by about
$1.43\times0.005=0.0071$ (exact: $0.0073$). At this operating point, one point of empirical risk is worth about $9.3$ nats of
KL: it is more profitable to improve the fit. The balance changes with $m$: for $m=100$, the same $5$ nats would be
worth $0.071$, five times more than the gain in empirical risk. This is exactly the trade-off that the Gibbs
posterior or a gradient descent on the certificate resolves automatically.

\exo{6.4}
\statement{$(\star)$ Derive the ``quadratic'' relaxation of the table of \cref{P-sec:sb_optim} from the refined Pinsker
inequality. Writing $a=\sqrt{\varepsilon/2}$, show that it is smaller than McAllester's bound iff $2a+2\sqrt{\hat q+a^2}\leq1$; in
particular, for small $\varepsilon$ it is the case as soon as $\hat q<\frac14$.}

\sol Let $p=\klup(\hat q,\varepsilon)>\hat q$. By the refined Pinsker inequality, $\varepsilon\geq\kl(\hat q\|p)\geq\frac{(p-\hat q)^2}{2p}$, so
$(p-\hat q)^2\leq2\varepsilon p$, i.e.\ $p^2-2(\hat q+\varepsilon)p+\hat q^2\leq0$. The roots of this quadratic are
$\hat q+\varepsilon\pm\sqrt{2\hat q\varepsilon+\varepsilon^2}$, so
\[
p\leq\hat q+\varepsilon+\sqrt{2\hat q\varepsilon+\varepsilon^2}=\Big(\sqrt{\hat q+\tfrac\varepsilon2}+\sqrt{\tfrac\varepsilon2}\Big)^2,
\]
where the last equality is checked by expanding: $(\sqrt{\hat q+a^2}+a)^2=\hat q+2a^2+2a\sqrt{\hat q+a^2}$ with $2a^2=\varepsilon$ and
$2a\sqrt{\hat q+a^2}=\sqrt{2\varepsilon}\sqrt{\hat q+\varepsilon/2}=\sqrt{2\hat q\varepsilon+\varepsilon^2}$. McAllester's bound is $\hat q+a$. Then
\[
(\sqrt{\hat q+a^2}+a)^2\leq\hat q+a\iff2a^2+2a\sqrt{\hat q+a^2}\leq a\iff2a+2\sqrt{\hat q+a^2}\leq1
\]
(dividing by $a>0$). When $\varepsilon\to0$, $a\to0$ and the condition becomes $2\sqrt{\hat q}<1$, i.e.\ $\hat q<\frac14$. So the quadratic
bound improves on McAllester's for all classifiers better than $25\%$ error, provided the complexity term is small; for
$\hat q$ close to $\frac12$, Pinsker's inequality is tight and McAllester's bound is better than this relaxation (but never better
than the kl itself).

%% ======================================================================
\section*{Section 7 -- Self-bounding algorithms for ensembles and beyond}
\addcontentsline{toc}{section}{Section 7 -- Self-bounding algorithms for ensembles and beyond}

\exo{7.1}
\statement{Derive the gradient \eqref{P-eq:grad_tnd_lambda} of the tandem objective, and check it numerically with finite
differences (notebook~4).}

\sol Write $c_1=\frac{4}{1-\lambda/2}$ and $c_2=\frac{4}{\lambda(1-\lambda/2)n_0}$, so that
$\mathcal B_{\mathrm{TND}}(\rho)=c_1\,\rho^\top\widehat{\mathbf L}\rho+c_2\big(2\KL(\rho\|\pi)+\ln\frac{2\sqrt{n_0}}\delta\big)$. The matrix $\widehat{\mathbf L}$ of
tandem losses is symmetric ($\widehat L_{tt'}=\widehat L_{t't}$), so $\nabla_\rho(\rho^\top\widehat{\mathbf L}\rho)=2\widehat{\mathbf L}\rho$. For the KL,
$\KL(\rho\|\pi)=\sum_t\rho_t\ln\rho_t-\rho_t\ln\pi_t$, and $\partial_{\rho_t}=\ln\rho_t+1-\ln\pi_t$, i.e.\
$\nabla_\rho\KL=\ln\frac\rho\pi+\mathbf1$ (componentwise). Hence
\[
\nabla_\rho\mathcal B_{\mathrm{TND}}=4\Big[\frac{2\widehat{\mathbf L}\rho}{1-\lambda/2}+\frac{2(\ln\frac\rho\pi+\mathbf 1)}{\lambda(1-\lambda/2)n_0}\Big].
\]
With the softmax parametrization $\rho=\mathrm{softmax}(\theta)$, the Jacobian is $\frac{\partial\rho_s}{\partial\theta_t}=\rho_s(\ind{s=t}-\rho_t)$,
so $\nabla_\theta\mathcal B=\rho\odot(g-(\rho^\top g)\mathbf 1)$ with $g=\nabla_\rho\mathcal B$, which is \eqref{P-eq:softmax_grad}. (The constant $\mathbf1$
in the KL gradient disappears in this projection, as it should: it corresponds to the direction that changes
$\sum_t\rho_t$.) \emph{Numerical check.} On the forest of the running example ($n_0=423$), at random logits $\theta$ and
$\lambda=0.7$, the centred finite differences with step $10^{-6}$ agree with the analytical gradient up to a relative error
of $6\times10^{-8}$:
\begin{verbatim}
num = np.array([(f(th + 1e-6*e) - f(th - 1e-6*e)) / 2e-6
                for e in np.eye(100)])
print(np.linalg.norm(g - num) / np.linalg.norm(num))   # 5.6e-08
\end{verbatim}

\exo{7.2}
\statement{Explain why the uniform prior over the trees of a random forest would \emph{not} be a valid prior if the empirical
losses were computed on the whole training sample.}

\sol The uniform prior on the \emph{indices} of the trees is fixed in advance, and in this sense it is legitimate.
The problem lies in the third step of the proof of Theorem~\ref{P-thm:general}: for each index $t$, we need the empirical loss
$\widehat L_t$ to be an average of i.i.d.\ variables of mean $\RD(h_t)$, independent of $h_t$. If $\widehat L_t$ is computed on the whole
training sample, it includes the examples used to grow $h_t$, on which a fully grown tree makes almost no error: $\widehat L_t$ is
then a strongly biased estimate of $\RD(h_t)$ (close to $0$ for unpruned trees), and Maurer's lemma does not apply.
Equivalently, the prior on the trees $h_t$ themselves (not on their indices) depends on the sample used to evaluate
them. The certificate would be meaningless: it would certify a near-zero risk for a forest whose trees err on $19\%$ of
the test digits. The out-of-bag construction restores the independence by evaluating each tree only on the examples
it has not seen (Theorem~\ref{P-thm:oob}).

\exo{7.3}
\statement{In PBGD, show that when all the training points have a large normalized margin, the gradient of the loss term
vanishes, and interpret the resulting value of $\|\mathbf w\|$.}

\sol The gradient of the objective is $\nabla F(\mathbf w)=-C\sum_i\varphi(t_i)\frac{y_ix_i}{\|x_i\|}+\mathbf w$ with
$t_i=\frac{y_i\langle\mathbf w,x_i\rangle}{\|x_i\|}$ and $\varphi$ the standard Gaussian density. Since $\|\frac{y_ix_i}{\|x_i\|}\|=1$, the loss
part has norm at most $C\sum_i\varphi(t_i)\leq Cm\,\varphi(t_{\min})$, and $\varphi(t)=\frac{e^{-t^2/2}}{\sqrt{2\pi}}\to0$ very fast: if all
$t_i\geq t_{\min}$ large, the loss gradient vanishes. Then $\nabla F(\mathbf w)\approx\mathbf w\neq0$, and gradient descent shrinks
$\mathbf w$. Shrinking $\mathbf w$ decreases all normalized margins proportionally ($t_i$ is linear in $\|\mathbf w\|$), until the smallest
ones become of order $1$--$2$, where $\varphi(t_i)$ is no longer negligible. At a stationary point,
$\mathbf w=C\sum_i\varphi(t_i)\frac{y_ix_i}{\|x_i\|}$: $\mathbf w$ is a combination of the points with moderate margins, the analogue
of support vectors. If the data are separable with geometric normalized margin $\gamma$ in direction $\mathbf u$, the
stationary $\mathbf w=s\mathbf u$ has $s\gamma=O(1)$, i.e.\ $\|\mathbf w\|\approx\frac{c}{\gamma}$ with $c$ growing slowly (logarithmically) with $C$, since $\varphi(t)$ decays like
$e^{-t^2/2}$. The KL, $\frac12\|\mathbf w\|^2\approx\frac{c^2}{2\gamma^2}$, is thus inversely proportional to the squared margin: large margins
give small certificates, which is the PAC-Bayesian explanation of the SVM. On the example of \cref{P-fig:pbgd}, the
norm grows from $3.5$ ($C=1$) to $13.6$ ($C=100$) only because of the noisy points, which pull $\mathbf w$ as long as their
margin is moderate.

\exo{7.4}
\statement{With notebook~4, compute the test disagreement of the C-bound posterior of \cref{P-tab:rf_posteriors} and check
whether the condition $d\geq R$ of Proposition~\ref{P-prop:compare} holds. Why is its certificate nevertheless larger than the
tandem certificate?}

\sol For the posterior minimizing the PAC-Bayesian C-bound, the test Gibbs risk is $R=0.176$ and the test
disagreement is $d=0.246$: the condition $d\geq R$ holds, and indeed the \emph{oracle} C-bound is $0.171$, smaller than the
oracle tandem bound $4e=0.210$. Its certificate, however, is $0.531$, against $0.361$ for the TND-optimized posterior. The
reason is statistical. The certificate replaces $R$ by an upper confidence bound and $d$ by a lower confidence bound,
each with $\delta/2$: from the OOB estimates $\hat R=0.153$ ($n_1=747$) and $\hat d=0.224$ ($n_3=423$ pairs), with $\KL=0.73$,
we get $\overline r=0.212$ and $\underline d=0.147$. The disagreement is estimated on the smaller pairwise sample with a doubled KL,
and it enters through $\frac1{1-2d}$, which amplifies its estimation error: plugging $\overline r$ and $\underline d$ gives
$1-\frac{(1-0.424)^2}{1-0.294}=0.531$. The tandem certificate needs a single kl inversion of a small quantity ($\hat e=0.041$),
where the kl is at its best. This is the general lesson of Proposition~\ref{P-prop:compare}: the best oracle bound is not
necessarily the best certificate; the number of quantities to estimate, and their size, matter.

\exo{7.5}
\statement{$(\star)$ For a stochastic network with $D=10^5$ weights, $\sigma_j=\sigma_0=0.03$ and $\|\boldsymbol\mu-\boldsymbol\mu_0\|=3$, compute
$\KL(Q\|P)$. How many examples are needed for the complexity term $\KL/m$ to be below $0.01$? What does this suggest
about data-dependent priors?}

\sol With $\sigma_j=\sigma_0$, the terms $\frac{\sigma_j^2}{\sigma_0^2}-1+\ln\frac{\sigma_0^2}{\sigma_j^2}$ vanish, and
\[
\KL(Q\|P)=\frac12\sum_j\frac{(\mu_j-\mu_{0,j})^2}{\sigma_0^2}=\frac{\|\boldsymbol\mu-\boldsymbol\mu_0\|^2}{2\sigma_0^2}=\frac{9}{2\times0.0009}=5000\text{ nats}.
\]
The complexity term $\KL/m$ is below $0.01$ only when $m\geq5000/0.01=500\,000$ examples, about eight times the size of
MNIST's training set, and this ignores the confidence term and the kl inversion. Note that the dimension $D$ does not
appear: what matters is how far, in units of $\sigma_0$, the posterior mean has travelled from the prior mean. Moving
$3$ units of Euclidean distance with a noise scale of $0.03$ is enormous. This suggests two levers: take $\sigma_0$ as large
as the network tolerates (chosen on a grid with a union bound), and above all bring $\boldsymbol\mu_0$ close to $\boldsymbol\mu$ with a
data-dependent prior: if $\boldsymbol\mu_0$ is a network trained on half of the data, the final network may for instance move by a
distance of $0.3$ instead of $3$, which divides the KL by $100$ ($50$ nats): the certificate then becomes informative
with a few tens of thousands of examples. This is the approach of \citet{dziugaite2018data} and \citet{perez2021tighter}
(Proposition~\ref{P-prop:ddprior}).

\bibliographystyle{plainnat}
\bibliography{../common/pacbayes}
\end{document}
